
\documentclass{article}
\usepackage{colortbl}
\usepackage{makecell}
\usepackage{multirow}
\usepackage{supertabular}

\begin{document}

\newcounter{utterance}

\centering \large Interaction Transcript for game `cladder', experiment `full\_v1.5\_default', episode 8201 with selfrefine{-}Qwen3.5{-}27B.
\vspace{24pt}

{ \footnotesize  \setcounter{utterance}{1}
\setlength{\tabcolsep}{0pt}
\begin{supertabular}{c@{$\;$}|p{.15\linewidth}@{}p{.15\linewidth}p{.15\linewidth}p{.15\linewidth}p{.15\linewidth}p{.15\linewidth}}
    \# & \multicolumn{2}{c}{Player} && \multicolumn{2}{c}{Game Master} \\
    \hline

    \theutterance \stepcounter{utterance}  
    & & & \multicolumn{4}{p{0.6\linewidth}}{
        \cellcolor[rgb]{0.9,0.9,0.9}{
            \makecell[{{p{\linewidth}}}]{
                \texttt{\tiny{[P1$\langle$GM]}}
                \texttt{You are an expert in causal inference. The following question is not a typical commonsense query, but rather a meticulously designed question created by a professor specializing in causal inference, intended to assess the students' mastery of the course content.} \\
\texttt{Imagine a self{-}contained, hypothetical world with only the following conditions, and without any unmentioned factors or causal relationships: Tijv has a direct effect on tijw and xevo. Xevo has a direct effect on gyzp. Tijw has a direct effect on gyzp. For those who are not xevo and are not tijw, the probability of gyzp is 76\%. For those who are not xevo and are tijw, the probability of gyzp is 76\%. For those who are xevo and are not tijw, the probability of gyzp is 12\%. For those who are xevo and are tijw, the probability of gyzp is 14\%. For those who are not xevo, the probability of tijw is 28\%. For those who are xevo, the probability of tijw is 14\%. For those who are xevo, would it be more likely to see gyzp if the individual was not xevo?} \\
\texttt{Start your answer with 'Yes' or 'No', followed by additional reasoning or evidence to support your explanation.} \\
            }
        }
    }
     \\ \\

    \theutterance \stepcounter{utterance}  
    & & \multicolumn{4}{p{0.6\linewidth}}{
        \cellcolor[rgb]{0.95,0.95,0.95}{
            \makecell[{{p{\linewidth}}}]{
                \texttt{\tiny{[P1$\leftrightarrow$P]}}
                \texttt{{[}DRAFT{]}} \\
\texttt{No} \\
\\ 
\texttt{**Reasoning:**} \\
\\ 
\texttt{To answer this question, we must calculate the counterfactual probability of the outcome \$gyzp\$ for an individual who is currently \$xevo\$, under the hypothetical intervention where they are set to "not \$xevo\$" (\$do(\textbackslash{}neg xevo)\$).} \\
\\ 
\texttt{1.  **Identify the Causal Structure:**} \\
\texttt{    *   \$Tijv \textbackslash{}to tijw\$} \\
\texttt{    *   \$Tijv \textbackslash{}to xevo\$} \\
\texttt{    *   \$xevo \textbackslash{}to gyzp\$} \\
\texttt{    *   \$tijw \textbackslash{}to gyzp\$} \\
\texttt{    *   Note: \$xevo\$ and \$tijw\$ are confounded by \$Tijv\$ (they share a common cause), but there is no direct arrow between \$xevo\$ and \$tijw\$.} \\
\\ 
\texttt{2.  **Analyze the Intervention:**} \\
\texttt{    The question asks: "For those who are \$xevo\$, would it be more likely to see \$gyzp\$ if the individual was not \$xevo\$?"} \\
\texttt{    This requires comparing:} \\
\texttt{    *   \$P(gyzp | xevo)\$: The observed probability of \$gyzp\$ given the individual is \$xevo\$.} \\
\texttt{    *   \$P(gyzp | do(\textbackslash{}neg xevo))\$: The probability of \$gyzp\$ if we intervene to make the individual \$\textbackslash{}neg xevo\$.} \\
\\ 
\texttt{    *Crucial Causal Inference Point:* When we perform an intervention \$do(\textbackslash{}neg xevo)\$, we sever the incoming arrow to \$xevo\$ from \$Tijv\$. However, the distribution of the other variables in the system (specifically \$tijw\$) depends on whether they are affected by the intervention.} \\
\texttt{    *   \$tijw\$ is a child of \$Tijv\$.} \\
\texttt{    *   \$xevo\$ is a child of \$Tijv\$.} \\
\texttt{    *   In a standard causal model, intervening on \$xevo\$ does **not** change the distribution of its parent \$Tijv\$, nor does it change the distribution of its sibling \$tijw\$ (unless there is a direct link \$xevo \textbackslash{}to tijw\$, which the prompt explicitly excludes).} \\
\texttt{    *   Therefore, the distribution of \$tijw\$ under the intervention \$do(\textbackslash{}neg xevo)\$ remains the same as the marginal distribution of \$tijw\$ in the population, or more precisely, we must look at the specific sub{-}population defined by the counterfactual.} \\
\\ 
\texttt{    However, the phrasing "For those who are \$xevo\$..." usually implies a counterfactual query: \$P(gyzp\_\{\textbackslash{}neg xevo\} | xevo=1)\$.} \\
\texttt{    To calculate this, we need the distribution of \$tijw\$ for the specific individuals who are currently \$xevo\$, but under the world where they are \$\textbackslash{}neg xevo\$.} \\
            }
        }
    }
    & \\ \\

    \theutterance \stepcounter{utterance}  
    & & \multicolumn{4}{p{0.6\linewidth}}{
        \cellcolor[rgb]{0.95,0.95,0.95}{
            \makecell[{{p{\linewidth}}}]{
                \texttt{\tiny{[P1$\leftrightarrow$P]}}
                \texttt{{[}SELF{-}FEEDBACK{]}} \\
\texttt{FORMAT: problem} \\
\texttt{RULES: ok} \\
\texttt{STRATEGY: problem} \\
\texttt{REASON: The answer fails to provide the final conclusion or complete the calculation, stopping abruptly after setting up the counterfactual framework without answering the specific question asked.} \\
            }
        }
    }
    & \\ \\

    \theutterance \stepcounter{utterance}  
    & \multicolumn{4}{p{0.6\linewidth}}{
        \cellcolor[rgb]{0.561,0.737,0.561}{
            \makecell[{{p{\linewidth}}}]{
                \texttt{\tiny{[P1$\rangle$GM]}}
                \texttt{No} \\
\\ 
\texttt{**Reasoning:**} \\
\\ 
\texttt{To answer this question, we must compare the observed probability of \$gyzp\$ for individuals who are \$xevo\$ against the counterfactual probability of \$gyzp\$ for those same individuals if they were not \$xevo\$.} \\
\\ 
\texttt{1.  **Calculate the Observed Probability \$P(gyzp | xevo)\$:**} \\
\texttt{    We are given the conditional probabilities of \$gyzp\$ based on the status of \$xevo\$ and \$tijw\$:} \\
\texttt{    *   \$P(gyzp | xevo, \textbackslash{}neg tijw) = 0.12\$} \\
\texttt{    *   \$P(gyzp | xevo, tijw) = 0.14\$} \\
\\ 
\texttt{    We are also given the probability of \$tijw\$ given \$xevo\$:} \\
\texttt{    *   \$P(tijw | xevo) = 0.14\$} \\
\texttt{    *   Therefore, \$P(\textbackslash{}neg tijw | xevo) = 1 {-} 0.14 = 0.86\$.} \\
\\ 
\texttt{    Using the law of total probability:} \\
\texttt{    \$\$P(gyzp | xevo) = P(gyzp | xevo, \textbackslash{}neg tijw)P(\textbackslash{}neg tijw | xevo) + P(gyzp | xevo, tijw)P(tijw | xevo)\$\$} \\
\texttt{    \$\$P(gyzp | xevo) = (0.12 \textbackslash{}times 0.86) + (0.14 \textbackslash{}times 0.14)\$\$} \\
\texttt{    \$\$P(gyzp | xevo) = 0.1032 + 0.0196 = 0.1228\$\$} \\
\\ 
\texttt{2.  **Calculate the Counterfactual Probability \$P(gyzp | do(\textbackslash{}neg xevo))\$ for the specific sub{-}population:**} \\
\texttt{    The question asks about individuals who *are* \$xevo\$ in the actual world, but asks what would happen if they were *not* \$xevo\$. This is a counterfactual query: \$P(gyzp\_\{\textbackslash{}neg xevo\} | xevo=1)\$.} \\
\\ 
\texttt{    In this causal structure (\$Tijv \textbackslash{}to xevo\$ and \$Tijv \textbackslash{}to tijw\$), \$xevo\$ and \$tijw\$ are confounded by \$Tijv\$. The fact that an individual is \$xevo\$ provides information about their latent \$Tijv\$ status, which in turn influences their \$tijw\$ status.} \\
\\ 
\texttt{    However, the prompt provides the conditional probabilities \$P(tijw | xevo)\$ and \$P(tijw | \textbackslash{}neg xevo)\$.} \\
\texttt{    *   \$P(tijw | xevo) = 0.14\$} \\
\texttt{    *   \$P(tijw | \textbackslash{}neg xevo) =} \\
            }
        }
    }
    & & \\ \\

    \theutterance \stepcounter{utterance}  
    & & & \multicolumn{2}{p{0.3\linewidth}}{
        \cellcolor[rgb]{0.9,0.9,0.9}{
            \makecell[{{p{\linewidth}}}]{
                \texttt{\tiny{[GM$|$GM]}}
                \texttt{No} \\
            }
        }
    }
    & & \\ \\

    \theutterance \stepcounter{utterance}  
    & & & \multicolumn{2}{p{0.3\linewidth}}{
        \cellcolor[rgb]{0.9,0.9,0.9}{
            \makecell[{{p{\linewidth}}}]{
                \texttt{\tiny{[GM$|$GM]}}
                \texttt{yes} \\
            }
        }
    }
    & & \\ \\

    \theutterance \stepcounter{utterance}  
    & & & \multicolumn{2}{p{0.3\linewidth}}{
        \cellcolor[rgb]{0.9,0.9,0.9}{
            \makecell[{{p{\linewidth}}}]{
                \texttt{\tiny{[GM$|$GM]}}
                \texttt{game\_result = LOSE} \\
            }
        }
    }
    & & \\ \\

\end{supertabular}
}

\end{document}
